\section{引言：从分散局部结果到共同问题}
孪生素数问题要求在固定移位 $2$ 上证明无穷多个素数对。它并不只要求建立一个形式上类似相关和的表达式，也不只要求某个有限矩阵具有小范数。若一种方法把原始计数问题转移到新的载体，那么载体上的系数、支撑、量词与归一化必须一并转移；忽略其中任一项，可能得到一个正确但不能用于原问题的命题。

本研究档案包含编号 TPC1--TPC418 的多轮局部探索，以及若干原型、缺源标记和后续未编号研究。其形成过程是问题驱动的：一个表达式的障碍引出新的表达方式，一次源对应核查又要求重新辨认实际对象。把这些工作按编号顺次压缩，会保留探索顺序，却不一定保留数学逻辑。因此，本文以图~\ref{fig:map}为共同坐标，重建一条可独立阅读的论证，并把没有进入证明链的材料留在可追溯的补充索引中。

\begin{figure}[p]
\centering
\includegraphics[width=\textwidth]{../../../research/tpc-big-road/figures/tpc_route_cn.png}
\caption{用户提供的七岛／两桥研究地图，原样复用。岛号表示研究对象与任务分工，不是已完成的步骤。实线为拟议主路线，虚线为辅助路线，红色断桥表示图示中的待证转移，不能解释为所有方法不可能。特别地，动力学岛所写的事件质量下界是该路线需要落实的条件，而非本文对指定算术轨道已经证明的事实。图像 Bridge A 内含档案中的 Gate A 与 Gate B；图像 Bridge B 是指定种子动力学桥，与文件名中的 Bridge-B 不同。小字宜放大查看；配套 README 提供全分辨率原图入口。}
\label{fig:map}
\end{figure}

\subsection{一个框架，两种实质缺口}
解析路线从局部相关结构走向完整的有符号集体消去。本文固定讨论的一个目标量具有自然尺度 $xQ^2$；要成为既定端点契约的有效输入，需要严格超过 $1/400$ 的幂次节省。完整恒等式可以保证对象不被改变，但不会凭空提供这种节省。第三、四岛分别处理指定局部原子和实际出现项的合法对应，第二岛才承担相应估计。

动力学路线的困难不同。即使一个编码系统具有足够的事件质量和协方差控制，也仍须证明被原始算术问题指定的那条轨道位于良好点集。关于几乎所有初始点的结论，不包含任意预先指定点。这一量词缺口不能通过改称“典型轨道”消除。几何提升若只改变表示方式，也不能替代该输入。

\subsection{本文保留的核心结果}
本文围绕四组可以逐项检查的内容展开：
\begin{enumerate}
\item 在明确的固定移位系数、单位剩余类和核约定下，证明实际算子的精确分解与行能量规模，而不把模型矩阵当作原始算子。
\item 证明常向量范数下界、实际系数的范数下界与增长维数的低秩修补障碍，明确它们限制的是哪些上界方法，而不是原混合标量的大小。
\item 用完整 Kloosterman 恒等式及一短一完整数组的精确范数，分离代数表达、支撑压缩和有符号重组三项义务。
\item 构造一个具有统一导数控制的分割，并在原混合标量层面支付边界代价；把剩余内部估计和动力学种子回接列为未完成条件。
\end{enumerate}
动力学部分另给出实际移动筛事件的 $O(N)$ 方差界与 Haar 几乎处处无限命中证明；它与指定种子 $0$ 的未证结论分开陈述。这些是已有档案结果的统一陈述与证明重组，不在本文中宣称首次发现或最优性。早期有限证书与后续模型增长结果在附录中按证据类型区分，不被误计为实际算术收益。

\subsection{与既有数学输入的关系}
本文的素数计数输入是经典素数定理及其分部求和推论\cite{selberg1949}。混合系数的二阶矩估计使用 Mertens 乘积定理的无条件基线；文献\cite{lamzouri2015}的开头提供该基线的陈述，本文不使用其中关于偏差分布的条件性结论。频率分离、正半定压缩、有限 Fourier 正交和谱极小极大原理是结构证明中的标准工具，其使用不构成新的素数对定理。

Blomer--Pascadi 的临界区间双线性估计\cite{blomer2026}提供一种有用的外部基准：在正确长度与互素条件下，Kloosterman 双线性型具有局部幂次节省。本文只使用其定理~1.1在平方根长度处的专门化，并检查本研究的表达是否真的满足这类输入条件；不把外部定理的存在等同于完整物理表达的估计已经完成。

本文不是该领域文献的穷尽综述，也不以内部资料数量代替与外部文献的创新性比较。它是一份限定研究计划的数学综整母稿。正式投稿之前，仍须进行贡献定位、补充相关文献、核验选用的全部历史依赖，以及独立的全文数学审查。

\subsection{阅读路径}
第~\ref{sec:arithmetic}节给出共同目标和证据层次；第~\ref{sec:structure}节解释局部原子与实际载体；第~\ref{sec:operator}--\ref{sec:boundary}节是自足的解析核心。动力学与几何路线在第~\ref{sec:dynamics}节单独处理，两桥的条件账本在第~\ref{sec:bridges}节汇总。读者不需要依次阅读所有历史小稿才能理解正文；但精确出处、未选用材料及审核范围可在配套来源矩阵和全目录索引中查找。
